\documentclass[10pt]{article}
\usepackage[margin=0.9in]{geometry}
\usepackage[T1]{fontenc}
\usepackage{lmodern}
\usepackage{microtype}
\usepackage{booktabs}
\usepackage{enumitem}
\usepackage{xcolor}
\usepackage[hidelinks]{hyperref}
\setlist{nosep,leftmargin=1.4em}
\newcommand{\code}[1]{\texttt{\small #1}}

\title{\textbf{Phantom}\\[4pt]\Large Exact, Deterministic Testing of LLM-Serving KV Caches\\without GPUs or Model Weights}
\author{Sameer Ahmed\\\small\url{https://github.com/sameer-sde/phantom-kv}}
\date{November 2026}

\begin{document}
\maketitle

\begin{abstract}
LLM-serving engines share KV-cache memory between requests through paged
blocks, prefix caching, preemption, offloading to CPU memory or a separate
cache server, and prefill/decode disaggregation. Bugs in this bookkeeping
rarely crash: a request silently reads another request's KV, or the cache
silently stops being reused. They are hard to test because a legitimate
cache hit already perturbs floating-point outputs, because the dangerous
interleavings are rare, and because GPUs are expensive. We present
\textbf{Phantom}, a testing tool built on a weightless model that writes a
\emph{provenance tag} (writer, position, prefix chain hash) into every KV
slot instead of K/V vectors, and checks every slot it reads. Two slots are
interchangeable exactly when their chain hashes match, which is the rule
prefix caching is allowed to exploit, so the in-model check raises no false
alarms on correct sharing. Combined with seeded, step-by-step execution and
delta-debugging shrinking, Phantom runs inside real, unmodified vLLM on a
laptop CPU. It caught every class of bug we planted (in vLLM's block pool,
its connector API and LMCache's transfer code), found no bugs in vLLM's two
CPU offloading systems, and independently found two real, open bugs in
LMCache: a server restart shorter than one heartbeat is never noticed, and
aborting a request whose cache lookup has finished but which has not been
scheduled leaks read locks for 300\,s. Both were already reported upstream;
Phantom shrank the second to a 7-event reproduction and provided A/B
evidence that the proposed fixes work.
\end{abstract}

\section{Introduction}
Modern LLM servers obtain their throughput by sharing KV-cache memory.
PagedAttention~\cite{kwon2023vllm} splits the cache into fixed-size blocks;
prefix caching lets requests with a common prefix share blocks; preemption
evicts and later recomputes requests; offloading copies blocks to CPU memory
or to an external cache server such as LMCache~\cite{liu2025lmcache}; and
disaggregated serving~\cite{zhong2024distserve,patel2024splitwise} computes a
prompt's KV on one engine and streams it to another. Each mechanism is
bookkeeping about who owns which block, and each can fail silently.

Testing this bookkeeping is difficult for three reasons. First, there is no
clean oracle: a correct prefix-cache hit already changes floating-point
results slightly, so ``the output changed'' does not imply a bug. Second, the
dangerous cases, such as a client aborting while its KV is still being
written, depend on interleavings that rarely reproduce. Third, fuzzing on
GPUs is slow and costly.

Phantom addresses all three:
\begin{itemize}
\item \textbf{An exact oracle.} A weightless model writes provenance tags
instead of K/V vectors and verifies every slot it reads (\S\ref{sec:oracle}).
\item \textbf{Deterministic execution.} One seed fixes the workload, the
timing of aborts and transfers, and the injected faults; runs replay
exactly and failures shrink with ddmin~\cite{zeller2002ddmin}
(\S\ref{sec:dst}).
\item \textbf{Real engines, no GPU.} Phantom plugs into vLLM's CPU backend
as an out-of-tree model and drives the real engine one step at a time;
vLLM's and LMCache's source code is not modified (\S\ref{sec:real}).
\item \textbf{Evidence of coverage.} Every run reports the engine's own
counters (prefix hits, preemptions, transfers, offloaded blocks), so a clean
result is backed by proof that the stressed path ran.
\end{itemize}
We evaluate Phantom against planted bugs and against code in which nobody
planted anything (\S\ref{sec:eval}), and describe two real LMCache bugs it
found (\S\ref{sec:cases}).

\section{Background}
In vLLM~\cite{kwon2023vllm}, a scheduler decides each step which requests run
and how many tokens each computes; a KV-cache manager assigns blocks from a
pool with reference counts and an LRU list of evictable cached blocks; a
prefix cache maps the hash of each full block's prefix to a block. A
\emph{KV connector} API lets external code load KV into a request's blocks
before it runs and save KV after it runs; it is how vLLM performs
disaggregated serving and offloading. vLLM ships two CPU offloading systems,
\code{SimpleCPUOffloadConnector} and the pluggable \code{OffloadingConnector}.
LMCache~\cite{liu2025lmcache} is a separate KV cache layer; in its
multi-process mode a cache server holds KV chunks (32 tokens here), and the
vLLM side looks chunks up asynchronously, read-locks hits until they are
loaded, and stores new chunks after compute.

\section{The provenance oracle}\label{sec:oracle}
Every request has a chain hash per position,
$\mathit{chain}[p] = H(\mathit{chain}[p-1], \mathit{token}[p])$, a 64-bit
BLAKE2b digest, so $\mathit{chain}[p]$ identifies the prefix
$\mathit{tokens}[0..p]$. The Phantom model writes
$\mathrm{Tag}(\mathit{writer}, p, \mathit{chain}[p])$ into the slot the
engine assigns to position $p$. When the model reads its context back
through the engine's block table, each slot is classified:
\begin{center}
\begin{tabular}{ll}
\toprule
Slot holds & Verdict \\
\midrule
same chain hash, any writer & correct (legal prefix sharing) \\
nothing & \code{unwritten-read} \\
a tag for a different position or chain & \code{misplaced-block} / \code{chain-break} \\
different chain hash, another request & \code{cross-request-leak} \\
\bottomrule
\end{tabular}
\end{center}
Two slots are interchangeable exactly when their chain hashes match, which
is precisely the condition under which prefix caching may share them. A
correct engine therefore never triggers a violation, up to a 64-bit hash
collision. The next token is a hash of everything read, so a bad read also
changes the visible output.

A second oracle compares every output, including partial outputs of aborted
requests, with a reference computed from the true prompt. The two are
complementary: a transfer that delivers \emph{another prompt's} complete KV
forms a valid chain for the wrong prefix and passes the in-model check, but
fails the output check; an empty slot fails the in-model check on its first
read. Phantom tests bookkeeping, not numerics: it cannot detect a wrong
attention kernel.

\section{Deterministic simulation and shrinking}\label{sec:dst}
Following deterministic simulation testing in distributed
systems~\cite{zhou2021fdb}, one seed is expanded into an explicit event
list: arrivals with prompt lengths and shared prefixes, client aborts at
chosen steps, and (for disaggregation) transfer delays. Delays are derived
from a hash of (seed, request, block), so removing events does not perturb
the timing of the rest. Because a run is a function of its event list,
failures replay exactly and can be minimized with ddmin~\cite{zeller2002ddmin}:
candidate event lists are re-run until no single event can be removed.
Phantom first ran against a small simulated engine that we wrote to
validate the oracle before trusting it on real code.

\section{Inside real engines}\label{sec:real}
\textbf{vLLM.} Phantom is registered through vLLM's plugin entry point as a
model architecture for the CPU backend (vLLM 0.33 development build). Two
details mattered. vLLM normally runs its engine core in a separate,
busy-looping process, so an abort scheduled for step 40 can land a step
late; Phantom runs the engine in-process and calls \code{step()} itself,
stepping while the scheduler reports pending work, as vLLM's own loop does.
And because vLLM caches blocks at schedule time, the model writes the whole
batch before checking, like real attention; an early version that did not
produced Phantom's only false alarm, which we traced and fixed.

\textbf{Faults and plants.} Planted bugs are installed at runtime by
wrapping engine functions, never by editing source: a stale prefix-cache
index entry (\code{BlockPool}), a refcount error, early block freeing on
abort during an asynchronous transfer, and transfer faults (another
request's KV, a load reported done but never written, a stale duplicate
copy landing later) in a connector loaded via vLLM's
\code{kv\_connector\_module\_path}.

\textbf{Offloading.} The CPU backend has no CUDA copy streams, so for both
offloading systems only the copy engine is replaced, with tensor copies
that follow the same instructions on separate load and store queues with
seeded delays; for \code{OffloadingConnector} this uses vLLM's own pluggable
\code{OffloadingSpec}. The scheduling and placement logic under test is
vLLM's.

\textbf{LMCache.} A CPU-only \code{lmcache server} (engine-driven transfers,
32-token chunks) runs beside the engine, with vLLM's own prefix cache off so
every reuse must come from LMCache. The harness waits for pending
asynchronous lookups before each step, clears the server between runs, can
restart it with a chosen downtime, and reads object and lock counts from its
\code{/status} endpoint.

\section{Evaluation}\label{sec:eval}
All runs used one Apple-silicon laptop CPU. CI re-runs the vLLM checks on
every pull request.

\textbf{Oracle validation (simulator, 1,000 seeds, 40 requests each, 25\%
aborts).} A correct engine produced 0 failing seeds; a planted
abort-during-transfer free produced 545; a planted stale prefix-cache entry
produced 1,000. Throughput was about 100 seeds per second on one core, and
failing runs shrank from about 50 events to a median of 6 (max 13).

\textbf{Planted bugs in real vLLM.} Table~\ref{tab:vllm} summarizes. A
planted refcount bug never reached Phantom: vLLM's assertion
\code{ref\_cnt == 0} fired first. Bugs that make the bookkeeping
\emph{inconsistent} trip assertions; bugs that leave it \emph{consistent but
wrong}, like a stale index entry, trip nothing, and that is where Phantom is
needed.

\begin{table}[h]
\centering\small
\begin{tabular}{lll}
\toprule
Setting & Unmodified vLLM & Planted bug \\
\midrule
Single engine, 20 seeds $\times$ 4 profiles & 0 violations, replay-identical & stale prefix cache caught on 15--20/20 \\
Shrinking (stale prefix cache) & --- & 38 events $\to$ 13 (235 vLLM runs) \\
Two engines, 200 seeds & 0/200, replay-identical & abort-free race 115/200 \\
 & & swap 199/200, silent drop 199/200 \\
 & & duplicate late copy 184/200 \\
Shrinking (abort-free race) & --- & 24 events $\to$ 2 \\
\bottomrule
\end{tabular}
\caption{Planted bugs in real vLLM. The profiles are baseline, chunked
prefill, preemption and 64-token blocks; the two-engine campaign exercised
1,371 asynchronous loads and 995 aborts, 335 of them while KV was in
flight. The one seed that swap and silent drop missed transferred no tokens
that were later read.}
\label{tab:vllm}
\end{table}

\textbf{Unplanted code: vLLM offloading.} For \code{OffloadingConnector}, a
50-seed campaign over eight configurations (including chunked offload,
revisits, preemption, slow stores and synchronous loads) stored 21,181
blocks and loaded 2,243 back, all verified, with every seed
replay-identical. A planted ``store reported early'' fault was caught on
24/50 seeds, and injected storage failures (117 failed loads, 757 failed
stores) were recovered correctly. \code{SimpleCPUOffloadConnector} was clean
over seven profiles. We found no bug in either system; because planted
faults were caught under the same conditions, we read this as a real
negative result for these configurations.

\textbf{LMCache.} With vLLM's prefix cache off, 11,040 tokens were served
from LMCache across 10 seeds, all verified, and a planted rotation of
destination blocks inside LMCache's scatter routine was caught on 10/10
seeds (1,149 misplaced blocks, 99 wrong outputs). Two further findings are
described next.

\section{Two real LMCache bugs}\label{sec:cases}
\subsection{A fast server restart is never noticed}
With the server restarted between workload runs, stores per server lifetime
were 19, 20, 21, 6 for a 25\,s outage but 19, 0, 0, 0 when the server
returned within a second: vLLM never re-registered, and caching stayed off
while outputs stayed correct. The worker re-registers only after a heartbeat
ping \emph{fails}, and the server answered pings positively even for a
worker it no longer knew, so a restart shorter than the 10\,s heartbeat
interval is invisible. The bug was already reported (LMCache issue \#4642)
with an open fix (PR \#6106) that makes the ping report registration. With a
fast restart followed by a full heartbeat, stores per lifetime went from 19,
0, 0, 0 without the fix to 19, 20, 21, 19 with it. On the fix's older base
commit the same restart \emph{hung} vLLM instead: the server rejected a
store from the forgotten worker but, in that version, sent no reply. Every
trial therefore ran with a time limit; a hang is a result.

\subsection{Read locks leaked by aborted requests}
Phantom's determinism check runs each seed twice and expects identical runs.
With LMCache, three of ten seeds differed on replay, although no output was
wrong. LMCache's \code{/status} showed objects still read-locked after every
request had finished, 1--5 per seed; a normal clear skips locked objects, so
they survived into the replay. With client aborts disabled, no locks
remained and every seed replayed identically. A single aborted request,
aborted at any state, never leaked, so we reused the ddmin shrinker with a
new failure predicate, ``this run left read locks'', and a forced clear
between candidates. It reduced a 41-event seed to 7 events in 171 runs,
leaking 3 locks on 5 of 5 re-runs. Logging the connector at the abort showed
the mechanism: the request was still waiting, its lookup had returned 96
tokens and read-locked 3 chunks, and the scheduler had not yet consumed the
result. The abort path dropped the lookup on the vLLM side without asking the
server to release the locks, which then expire only with the 300\,s read
TTL. This was reported as issue \#5656 with an open fix (PR \#6540). On the
fix's base commit the two minimal cases leaked 2 and 3 locks and 3 seeds
differed on replay; with the fix both were 0.

We posted both A/B results on the fix PRs rather than filing duplicates.

\section{Lessons}
\begin{itemize}
\item \textbf{Determinism is an oracle.} The lock leak never corrupted a
token; it surfaced only because a replay was not identical.
\item \textbf{A shrinker with a custom predicate is a debugger.} The 7-event
case was found, not guessed.
\item \textbf{Measure that the path ran.} A clean run with no prefix hits,
transfers or offloads tests nothing; counters made clean results meaningful
and exposed two harness mistakes of our own.
\item \textbf{Search before filing.} Both LMCache bugs were known; the useful
contribution was independent evidence about the proposed fixes.
\end{itemize}

\section{Limitations}
Everything ran on vLLM's CPU backend and a CPU-only LMCache server, so
GPU-only paths (CUDA copy streams, NIXL/RDMA transfers) are not covered, and
offloading was tested with a replacement copy engine. SGLang~\cite{zheng2024sglang}
and other engines are not covered yet. Phantom checks where KV goes, not
how it is computed. Workloads are small (tens of requests per seed) and
models tiny, which favors reaching edge cases quickly but may miss bugs that
need scale.

\section{Related work}
Deterministic simulation testing was popularized by
FoundationDB~\cite{zhou2021fdb} and is used by systems such as TigerBeetle;
Jepsen~\cite{jepsen} tests distributed databases with fault injection and
consistency checking. Property-based testing~\cite{claessen2000quickcheck}
and delta debugging~\cite{zeller2002ddmin} supply random generation and
shrinking. Phantom applies these ideas to LLM-serving KV caches, where its
contribution is an exact oracle obtained by replacing the model.
PagedAttention~\cite{kwon2023vllm}, SGLang~\cite{zheng2024sglang},
LMCache~\cite{liu2025lmcache} and disaggregated
serving~\cite{zhong2024distserve,patel2024splitwise} are the systems whose
memory sharing Phantom tests.

\section{Conclusion}
Replacing the model with one that writes provenance tags turns KV-cache
bookkeeping into something that can be checked exactly, on a laptop, inside
real engines. Phantom caught every planted bug class, gave a validated
negative result for vLLM's offloading, and independently found two real
LMCache bugs, one shrunk to seven events. Phantom is open source (MIT) and
installable with \code{pip install phantom-kv}; the repository contains the
harness, the minimal reproductions and a day-by-day development log.

\paragraph{Acknowledgments.} I used an AI coding assistant for code and log
analysis; I ran every experiment and checked every result.

\begin{thebibliography}{10}\small
\bibitem{kwon2023vllm} W.~Kwon, Z.~Li, S.~Zhuang, Y.~Sheng, L.~Zheng, C.~H.~Yu,
J.~E.~Gonzalez, H.~Zhang, I.~Stoica. Efficient Memory Management for Large
Language Model Serving with PagedAttention. \emph{SOSP}, 2023.
\bibitem{liu2025lmcache} Y.~Liu, J.~Yao, Y.~Cheng, Y.~An, X.~Chen, S.~Feng,
Y.~Huang, S.~Shen, R.~Zhang, K.~Du, J.~Jiang. LMCache: An Efficient KV Cache
Layer for Enterprise-Scale LLM Inference. arXiv:2510.09665, 2025.
\bibitem{zhou2021fdb} J.~Zhou, M.~Xu, A.~Shraer, B.~Namasivayam, A.~Miller,
E.~Tschannen, S.~Atherton, A.~J.~Beamon, et al. FoundationDB: A Distributed
Unbundled Transactional Key Value Store. \emph{SIGMOD}, 2021.
\bibitem{zeller2002ddmin} A.~Zeller, R.~Hildebrandt. Simplifying and Isolating
Failure-Inducing Input. \emph{IEEE Transactions on Software Engineering},
28(2), 2002.
\bibitem{claessen2000quickcheck} K.~Claessen, J.~Hughes. QuickCheck: A
Lightweight Tool for Random Testing of Haskell Programs. \emph{ICFP}, 2000.
\bibitem{jepsen} K.~Kingsbury. Jepsen. \url{https://jepsen.io}.
\bibitem{zheng2024sglang} L.~Zheng et al. SGLang: Efficient Execution of
Structured Language Model Programs. \emph{NeurIPS}, 2024.
\bibitem{zhong2024distserve} Y.~Zhong et al. DistServe: Disaggregating
Prefill and Decoding for Goodput-optimized Large Language Model Serving.
\emph{OSDI}, 2024.
\bibitem{patel2024splitwise} P.~Patel et al. Splitwise: Efficient Generative
LLM Inference Using Phase Splitting. \emph{ISCA}, 2024.
\end{thebibliography}
\end{document}
